\section{Separated coefficients and a denominator-free local construction}
\label{sec:diagonal-local}

The constructions in this section explain how an already known local
quadratic identity can be converted into finitely many fixed testing
features. They do not supply the missing global atlas of the quantum set.
The distinction is a quantifier issue: a construction on one rectangle
must still be extended to, or combined with a finite cover of, every
source and every testing point.

\subsection{An explicit separation identity for a positive coefficient}

Suppose
\[
 \gamma(x,y)=\sum_{i=1}^{n}a_i(x)b_i(y),\qquad
 \alpha_i=a_i(x_0),\quad \beta_i=b_i(y_0),\qquad
 c=\sum_i\alpha_i\beta_i>0.
\]
Write $a_i=\alpha_i+A_i$ and $b_i=\beta_i+B_i$. Choose a positive
rational number $\delta$ such that
\begin{equation}
 3n\delta^2+2\delta
 \max\left\{\sum_i|\alpha_i|,\sum_i|\beta_i|\right\}<c.
 \label{eq:diagonal-delta}
\end{equation}
Such a choice exists because the left side tends to zero with $\delta$.
Work on a product neighborhood where
$|A_i(x)|\le\delta$ and $|B_i(y)|\le\delta$ for every $i$. Define
\begin{align}
 p(x)&=\frac{c-n\delta^2}{2}
       +\sum_i(\beta_i-\delta)A_i(x),\nonumber\\
 q(y)&=\frac{c-n\delta^2}{2}
       +\sum_i(\alpha_i-\delta)B_i(y).
 \label{eq:diagonal-pq}
\end{align}
Then the following identity is exact:
\begin{equation}
 \boxed{\gamma(x,y)=p(x)+q(y)
   +\sum_i\bigl(A_i(x)+\delta\bigr)
              \bigl(B_i(y)+\delta\bigr).}
 \label{eq:diagonal-separated}
\end{equation}
Indeed, the $n\delta^2$ terms cancel those in $p+q$, the terms linear
in $A_i$ and $B_i$ reduce to $\beta_iA_i$ and $\alpha_iB_i$, and the
remaining terms are $A_iB_i$. Moreover,
\begin{align*}
 p(x)&\ge \frac c2-\frac{3n\delta^2}{2}
                      -\delta\sum_i|\beta_i|>0,\\
 q(y)&\ge \frac c2-\frac{3n\delta^2}{2}
                      -\delta\sum_i|\alpha_i|>0.
\end{align*}
The other factors in \eqref{eq:diagonal-separated} are nonnegative.
Thus a positive coefficient in an algebraic tensor product is, locally,
an explicit sum of $n+2$ products of nonnegative one-factor functions.
No inverse involving both factors is introduced.

For example, on $|x|,|y|\le1/4$,
\begin{equation}
 1+xy=(x+1/4)(y+1/4)
      +(15/32-x/4)+(15/32-y/4).
 \label{eq:diagonal-toy}
\end{equation}
The last two factors are at least $13/32$.

\paragraph{Conversion into fixed blocks of order two.}
For any separated nonnegative coefficient $a(x)b(y)$ and any fixed
scalar function $\ell$, put
\[
 v(y)=\sqrt{b(y)}\binom{1}{\ell(y)},\qquad
 Q(x)=a(x)
 \begin{pmatrix}\ell(x)^2&-\ell(x)\\-\ell(x)&1\end{pmatrix}.
\]
Then $Q(x)\succeq0$ and
\begin{equation}
 v(y)^TQ(x)v(y)=a(x)b(y)
                \bigl(\ell(y)-\ell(x)\bigr)^2.
 \label{eq:diagonal-block}
\end{equation}
The two entries of $v$ are fixed functions of the testing variable.
Applying this construction to the $n+2$ terms in
\eqref{eq:diagonal-separated} gives $n+2$ blocks of order two for
$\gamma(x,y)(\ell(y)-\ell(x))^2$.
If the factors are evaluated in independently labelled real closed
fields, their square roots are taken in their respective fields.
Equation \eqref{eq:diagonal-block} is therefore also an identity in
the corresponding tensor product. It uses no identification of the
two field copies.

\subsection{The smooth-diagonal statement and its hypotheses}

Let $Y\subset\mathbb R^m$ be a smooth affine real algebraic set near a
point $p$, with coordinate ring $B$ and ambient coordinate column $z$.
Write
\[
 \Delta z=z\otimes1-1\otimes z,\qquad
 J=\ker(B\otimes_{\mathbb R}B\longrightarrow B).
\]
The homomorphism is restriction to the diagonal. Suppose an exact
identity has already been supplied:
\begin{equation}
 \tau=\Delta z^T H\Delta z,\qquad
 H=H^T\in\operatorname{Mat}_m(B\otimes B),
 \label{eq:diagonal-H}
\end{equation}
and suppose the restriction of $H(p,p)$ to $T_pY$ is positive
definite. Membership $\tau\in J^2$ gives a representation of the form
\eqref{eq:diagonal-H}, because $J$ is generated by the coordinate
differences; the positive tangential residue is a separate hypothesis.
It is not implied by membership in $J^2$.

Under these hypotheses, there is a product neighborhood of $(p,p)$
on which $\tau$ has a finite fixed-feature representation using only
blocks of order two. This representation can be constructed without
inverting a mixed function of the two points.

\paragraph{Step 1: remove only the normal degeneracy.}
Choose defining equations $g_1,\ldots,g_r$ whose differentials span
the conormal space at $p$. Telescopic polynomial divided differences
give a matrix $A\in\operatorname{Mat}_{r\times m}(B\otimes B)$ with
\[
 A\Delta z=0,\qquad
 A(p,p)=Dg(p),\qquad \ker A(p,p)=T_pY.
\]
For example, every polynomial difference $g(z)-g(z')$ is a sum of
the coordinate differences multiplied by polynomial divided
differences. In $B\otimes B$ that polynomial difference is zero.

There is a positive integer $t$ for which
\[
 H_*=H(p,p)+tA(p,p)^TA(p,p)
\]
is positive definite on $\mathbb R^m$. To see this directly, decompose
$\mathbb R^m=T_pY\oplus(T_pY)^\perp$. The upper left block of
$H(p,p)$ is positive definite. On the normal summand,
$A(p,p)^TA(p,p)$ is positive definite. The Schur complement of the
upper left block becomes positive for all sufficiently large $t$.
Adding $tA^TA$ does not change \eqref{eq:diagonal-H}, since
$\Delta z^TA^TA\Delta z=0$ exactly.

\paragraph{Step 2: use a constant change of coordinates.}
Choose a real invertible constant matrix $C$ with $H_*=C^TC$ and put
\[
 f=Cz,\qquad
 \widehat H=C^{-T}(H+tA^TA)C^{-1}.
\]
Then $\widehat H(p,p)=I_m$ and
$\tau=\Delta f^T\widehat H\Delta f$. The inverses here are inverses
of a fixed real matrix, not inverses of functions on $Y\times Y$.

\paragraph{Step 3: use a fixed finite frame.}
Set $\varepsilon=1/(4m)$. For every symmetric matrix $K$ and vector
$d\in\mathbb R^m$, the polynomial identity
\begin{align}
 d^TKd={}&\sum_i\bigl(K_{ii}-2(m-1)\varepsilon\bigr)d_i^2
 \nonumber\\
 &+\sum_{i<j}\left(\varepsilon+\frac{K_{ij}}2\right)(d_i+d_j)^2
 \nonumber\\
 &+\sum_{i<j}\left(\varepsilon-\frac{K_{ij}}2\right)(d_i-d_j)^2
 \label{eq:diagonal-frame}
\end{align}
follows by expansion. Apply it with $K=\widehat H$ and
$d=\Delta f$. At $(p,p)$ its coefficients are
$(m+1)/(2m)$ on the $m$ coordinate squares and $1/(4m)$ on the
pair squares. All are strictly positive. They remain positive on a
sufficiently small product neighborhood.

There are $m^2$ fixed frame functions:
$f_i$ and $f_i+f_j,f_i-f_j$ for $i<j$. Each coefficient lies in
$B\otimes B$ and therefore has a finite separated expansion.
Apply \eqref{eq:diagonal-separated} to each coefficient, shrink the
two neighborhoods finitely many times, and use
\eqref{eq:diagonal-block}. This proves the claimed local construction.

\paragraph{What denominators are permitted.}
Pure-factor inverses can be handled by adjoining, to that factor,
a coordinate satisfying an equation such as $rh=1$, on the locus
where $h\ne0$. This is still a one-factor construction. The analogous
move for a mixed denominator need not have finite separation rank.
For instance,
\[
 K(x,y)=\frac1{1-xy}
\]
has infinite separation rank on any product of open intervals about
zero. If it were a sum of $N$ separated functions, choose $N+1$
distinct nonzero $x_j$ in the first interval. The functions
$1/(1-x_jy)$ would lie in an $N$-dimensional space. They are linearly
independent: a linear relation, multiplied by
$\prod_j(1-x_jy)$, is a polynomial identity, and evaluation at
$y=1/x_j$ kills all but the $j$th coefficient. This contradiction
does not require those poles to lie in the original interval.

The theorem above therefore requires the actual $J^2$ identity in
the indicated tensor product, not a representation obtained by an
uncontrolled mixed rational inversion. It also leaves the global
finite-cover question open.

\section{Exact extraction of the displayed \texorpdfstring{$B_0$}{B0} coefficient module}
\label{sec:b0-module}

The recovered August 30 source displays, in its equation (12.1), a
matrix with four scalar coefficients. The August 31 correction R5
specifies its factor typing: $A,T,p,w,h,\varepsilon_0$ belong to the
point factor, and $U,V$ belong to the support factor. We retain that
typing. For the algebraic calculation below, $c_i$ simply denote
the four coefficient-factor scalars; the phrase ``testing variables''
means $(U,V,z)$ in this module, without asserting a new physical
source-to-point identification.

Let $e_0,\ldots,e_3$ be the standard basis of $\mathbb R^4$ and
$\omega=(U,-1,0,-UV)^T$. Define
\begin{align}
 S_c(U,V)={}&c_1(1+U^2)e_0e_0^T
       +c_2(1+V^2)\omega\omega^T\nonumber\\
 &+c_3(1+U^2)(1+V^2)e_1e_1^T+c_4U^2e_2e_2^T.
 \label{eq:b0-module}
\end{align}
Then
\begin{equation}
 S_c(U,V)\succeq0\text{ for every real }U,V
 \quad\Longleftrightarrow\quad c_1,c_2,c_3,c_4\ge0.
 \label{eq:b0-orthant}
\end{equation}
The reverse implication follows from the displayed sum. The forward
implication follows from four rational tests:
\begin{center}
\begin{tabular}{ccl}
\toprule
$(U,V)$ & state $z$ & $z^TS_c(U,V)z$\\
\midrule
$(0,0)$ & $e_0$ & $c_1$\\
$(1,1)$ & $e_3$ & $2c_2$\\
$(1,1)$ & $(0,1,0,-1)^T$ & $4c_3$\\
$(1,0)$ & $e_2$ & $c_4$\\
\bottomrule
\end{tabular}
\end{center}
Thus this four-dimensional coefficient cone has the evaluated pencil
$\operatorname{diag}(c_1,c_2,c_3,c_4)\succeq0$, with order four and
no auxiliary variables.

For $z=(z_0,z_1,z_2,z_3)^T$, a fixed nine-function certificate is
\begin{align*}
 v={}&\bigl(z_0,\ Uz_0,\ Uz_0-z_1-UVz_3,\
 UVz_0-Vz_1-UV^2z_3,\\
 &\hspace{25mm} z_1,\ Uz_1,\ Vz_1,\ UVz_1,\ Uz_2\bigr)^T,
\end{align*}
and the exact identity is
\[
 z^TS_cz
 =v^T\operatorname{diag}(c_1,c_1,c_2,c_2,
                         c_3,c_3,c_3,c_3,c_4)v.
\]
For the original displayed coefficients,
\begin{align*}
 p&=1-A-LT^2,&w&=1-LT^2,&h&=1-L(1+T^2),\\
 c_1&=\frac{AT(1+T^2)p}{w},&
 c_2&=A\varepsilon_0,&c_3&=\varepsilon_0p,\\
 c_4&=\frac{A\varepsilon_0(1+T^2)wp}{h}.&&&
\end{align*}
At the rational strict-chamber point
$(T,L,A,\varepsilon_0)=(1,2/5,1/5,1)$, these become
\[
 (p,w,h)=(2/5,3/5,1/5),\qquad
 (c_1,c_2,c_3,c_4)=(4/15,1/5,2/5,12/25).
\]
The calculation also gives $L(T^2+2)-1=1/5>0$.

\paragraph{A precise conditional congruence envelope.}
Every displayed feature belongs to
\[
 \operatorname{span}\{1,U,V,UV,UV^2\}
   \otimes\operatorname{span}\{z_0,z_1,z_2,z_3\}.
\]
Consequently, after a $4\times4$ state congruence
$z=R\widetilde z$ whose entries depend only on the coefficient
factor, the twenty fixed functions
\[
 \{\widetilde z_i,U\widetilde z_i,V\widetilde z_i,
                    UV\widetilde z_i,UV^2\widetilde z_i:0\le i<4\}
\]
contain every transformed feature. A congruence depending arbitrarily
on $U,V$ is not covered by this conclusion. This containment also
does not assert that an order-two block bound survives an arbitrary
change of state coordinates.

The recovered text describes a pure-point radial congruence and a
common positive clearing before (12.1), but does not print those
maps. Therefore the exact module, its orthant, and the conditional
twenty-function envelope are complete; the restoration of this
module to the full original-coordinate slack remains an explicit
missing input.

\section{Reconstructed event-contact cofactors}
\label{sec:event-cofactors}

This calculation starts from the canonical matrix in equation (4.2)
of the recovered August 30 source. Restrict to the kernel of its
$\lambda$ row by writing
\[
 \psi=(-vy_3,y_3,y_1,y_2)^T,
\]
and set the two extra coefficients $\alpha=\beta=0$. The restricted
quadratic form has coefficients
\begin{align}
 G_d={}&2\bigl[(D_v+u^2v^2)y_1^2-2u^2vy_1y_2+u^2y_2^2
                  +D_vy_3^2+2uD_vy_2y_3\bigr],\nonumber\\
 G_1={}&4D_uD_vy_3^2,\qquad
 G_2=4(uvy_1-uy_2-D_vy_3)^2,\nonumber\\
 G_3={}&4D_v(y_1+uvy_3)^2,\qquad
 G_4=4D_uD_vy_1^2,\qquad
 G_5=4(y_1+vy_2)^2,
 \label{eq:event-restriction}
\end{align}
where $D_u=1+u^2$ and $D_v=1+v^2$. These six restrictions have been
checked directly against the canonical row vectors. Substituting
$u=u_1/u_0$, $v=v_1/v_0$ and multiplying the five positive $G_i$ by
$u_0^2v_0^4$ gives exactly four times the five homogeneous forms in
the recovered equation (11.0). This is an algebraic homogenization
identity, including its polynomial extension to projective endpoints.

On the state chart $y_2=1$, put $a=y_1$ and $b=y_3$. Let
$G=(G_d,G_1,\ldots,G_5)$ after this substitution and form
\[
 J=\begin{pmatrix}
 G\\ \partial_uG\\ \partial_vG\\ \partial_aG\\ \partial_bG
 \end{pmatrix},\qquad
 X_j=(-1)^j\det J_{\widehat j},\quad 0\le j\le5.
\]
Here $J_{\widehat j}$ is the five-column matrix obtained by deleting
column $j$, and indexing starts at zero. Laplace expansion gives
$JX=0$ identically. Direct exact calculation over
$\mathbb Q[u,v,a,b]$ gives the common factor
\[
 g=8192\,b(1+u^2)(1+v^2)^3.
\]
After division by $g$, the primitive polynomial vector
$\widetilde X=X/g$ has component degrees
\[
 (15,13,13,11,15,17)
\]
and respectively
\[
 (59,40,43,36,42,38)
\]
nonzero monomials. Its component gcd is one and
$J\widetilde X=0$ is checked as five exact polynomial identities.
All coefficients, the six defining forms, and the thirty Jacobian
entries are printed as machine-readable rational polynomial records
in \nolinkurl{diagonal/EVENT_CONTACT_EXACT.json}.

These cofactors parametrize stationary solutions of the restricted
contact equations wherever the contact matrix has rank five. They do
not by themselves prove componentwise nonnegativity, physical support
validity, a full four-state kernel equation, or completeness at rank
drop and projective boundaries. Those additional claims must be
proved separately before this vector can be used as a global source
parametrization.

\section{The finite-cover assembly that remains to be instantiated}
\label{sec:finite-cover-conditional}

The final passage from local fixed features to an affine pencil is
elementary once the necessary global data have actually been supplied.
For clarity, the statement below lists those data rather than
presupposing them.

Let $T$ be a testing set with a map $\Phi:T\to\mathbb R^d$, and let
$C=\operatorname{conv}\Phi(T)$ be closed and full dimensional. Let
$S$ index every affine functional nonnegative on $C$, including its
constant. Suppose that finitely many rectangles
$U_i\times V_i\subset S\times T$ cover all of $S\times T$. On each
rectangle, assume a supplied exact slack identity
\begin{equation}
 f_s(\Phi(t))=\sum_k v_{ik}(t)^TQ_{ik}(s)v_{ik}(t),\qquad
 Q_{ik}(s)\succeq0\quad(s\in U_i),
 \label{eq:finite-cover-local}
\end{equation}
with finitely many fixed vector functions $v_{ik}$ of order at most
two. All normalizations, coordinate changes, and clearing factors
must already have been restored in this identity.

\paragraph{Step 1: convert the rectangle cover into fixed functions.}
Partition $S$ by its membership patterns in the finitely many $U_i$,
and partition $T$ by its membership patterns in the $V_i$. If $A$ and
$B$ are nonempty cells of these partitions, the covering property
provides an $i=i(A,B)$ such that $A\times B\subset U_i\times V_i$.
This follows because membership in every $U_i,V_i$ is constant on a
cell. Choose one such index once for each pair of cells.

For each pair $(A,B)$, extend $v_{i(A,B),k}$ by zero outside $B$.
Extend its source matrix by zero outside $A$. For a fixed source
there is exactly one source cell, while every test lies in exactly
one test cell. Summing the resulting finitely many identities gives
a global representation with fixed testing functions and positive
source matrices. The indicator functions merely select cells; they
do not enlarge the vector order. If the original data and cells are
semialgebraic, these zero extensions are semialgebraic as well.

\paragraph{Step 2: form the finite product space.}
Let $v_j$ denote the resulting finite family of vector functions.
Define the real vector space
\[
 W=\operatorname{span}_{\mathbb R}
 \{1,\Phi_1,\ldots,\Phi_d,(v_j)_a(v_j)_b:j,a,b\}.
\]
Choose an actual basis
$w_0=1,w_1=\Phi_1,\ldots,w_d=\Phi_d,w_{d+1},\ldots,w_D$.
Full dimensionality ensures independence of the first $d+1$
functions. Record every product expansion
\[
 (v_j)_a(v_j)_b=\sum_{k=0}^{D}P_{j,ab,k}w_k.
\]
Only this finite linear span is needed; multiplication closure of
$W$ is not a requirement.

\paragraph{Step 3: write and prove the pencil.}
For $x\in\mathbb R^d$ and $y\in\mathbb R^{D-d}$, put
\begin{equation}
 M_j(x,y)_{ab}=P_{j,ab,0}
       +\sum_{\nu=1}^{d}x_\nu P_{j,ab,\nu}
       +\sum_{k=d+1}^{D}y_kP_{j,ab,k}.
 \label{eq:finite-cover-pencil}
\end{equation}
Then
\[
 C=\{x:\exists y,\ M_j(x,y)\succeq0\text{ for every }j\}.
\]
For the forward inclusion, if $x$ is a finite convex combination of
$\Phi(t)$, let $L$ be the same convex combination of evaluations on
$W$. Each $M_j=L(v_jv_j^T)$ is positive semidefinite, $L(1)=1$, and
$L(\Phi_\nu)=x_\nu$.

For the reverse inclusion, a feasible collection $(x,y)$ defines
the linear functional $L$ on the chosen basis of $W$. If $x\notin C$,
closed convex separation gives an affine $f_s$ nonnegative on $C$
with $f_s(x)<0$. Apply $L$ to its global identity assembled in
Step 1. The left side is $f_s(x)$, while the right side is a sum of
$\operatorname{tr}(Q_j(s)M_j(x,y))\ge0$, a contradiction.
Both inclusions are thus proved once the hypotheses and the actual
coefficient tables have been evaluated.

For the present full eight-coordinate quantum problem, those global
inputs have not all been supplied. The local constructions in this
release give usable pieces of the required data. They do not replace
the missing restoration maps, exceptional and boundary coverage, or
the finite all-source/all-test cover. Consequently the global feature
list, global evaluated pencil, and verification of both global
inclusions remain open.

\paragraph{Fresh verification records.}
The script \nolinkurl{diagonal/regenerate_diagonal.py} regenerates
the exact coefficient data and four algebraic verification groups
with SymPy 1.14.0. The corresponding receipt is
\nolinkurl{diagonal/DIAGONAL_VERIFICATION.json}. Its checks corroborate
the displayed identities and restricted cofactor computations. They
are not a proof-assistant verification of the analytical arguments
and do not certify the missing global hypotheses.
